% Part 2 chapter carrying the five Virial papers, in date order, read in full 2026-10-02 (2,370 extracted lines):
%  (1) Virial Efficiency and Effective Nonlinear Exponent (25 Feb 2026); (2) The Virial Partition Across a Wide Range of Physical Scales (25 Feb 2026);
%  (3) Dark Matter and Dark Energy as Virial Partners (Mar 2026); (4) Gravitational Decoherence, the Virial Partition, and the Emergence of Classical
%  Structure (Mar 2026); (5) The Thermodynamic Identity Governing the Virial Theorem (PRL version, 18 Mar 2026).
% Checks: docs/verification/virial/VIRIAL_CHECK.md, NBODY_TRACE.md. Interpretation sections of (3) and (4) go to Part 5 (marked below).
\providecommand{\authorhold}[1]{\par\smallskip\noindent\fbox{\parbox{0.95\linewidth}{\small\textbf{Author hold.} #1}}\par\smallskip}

\chapter{The virial partition: where $\beta_m=\Omega_m/2$ comes from}\label{ch:virial}

\section*{Overview}
Every result in this part of the book carries one number: the coupling $\beta_m=\Omega_m/2$. It is not fitted. It comes from the oldest equilibrium theorem in
mechanics, the virial theorem of Clausius (1870), which fixes the share of binding energy that a stable system bound by a $1/r$ potential holds as motion.
This chapter follows the five papers that built the case: the theorem itself and the reason its factor is exactly $1/2$; the evidence that the same $1/2$
holds from the hydrogen atom to galaxy clusters; the thermodynamic identity that tells us what the kinetic half is; the step from the theorem to the
cosmological coupling; and what the coupling then predicts for growth, for the two expansion rates, and for the tests now coming.

\section{The theorem and why the factor is one half}
For a bound system whose potential energy is a homogeneous function of degree $k$ in the coordinates, the time averages of kinetic and potential energy satisfy
\begin{equation}
2\langle K\rangle=k\,\langle V\rangle .
\end{equation}
This follows from Euler's theorem for homogeneous functions applied to the virial $\sum_i\mathbf r_i\cdot\mathbf p_i$, whose time average is constant for a
bound system. Gravity and the Coulomb interaction both have $V\propto r^{-1}$, so $k=-1$ and
\begin{equation}
2\langle K\rangle+\langle V\rangle=0,\qquad \langle K\rangle=\tfrac12|\langle V\rangle| .
\end{equation}
The factor is set by the degree of the potential, and the degree is set by geometry: in three spatial dimensions Gauss's law makes the field of a point source
fall as $1/r^2$, so the potential falls as $1/r$. Nothing in the result depends on the number of bodies, the mass distribution or the scale. The same
$1/2$ appears in the Bohr atom, in a star, and in a cluster of galaxies.

\section{The same one half across physical scales}
Table~\ref{tab:virial_domains} collects the domains in which the $1/r$ virial ratio has been tested.

\begin{table}[h]\centering\small
\caption{The virial ratio for $1/r$ binding, from atoms to clusters of galaxies.}\label{tab:virial_domains}
\begin{tabular}{llll}
\hline
System & Scale & Quantity & Result\\\hline
Hydrogen atom & $10^{-10}$ m & $\langle K\rangle/|\langle V\rangle|$ & $1/2$ exactly (13.6\,eV of 27.2\,eV)\\
20 atoms, H--Xe (Hartree--Fock limit) & $10^{-10}$ m & $T/|V|$ & $1/2$ (exact for converged wavefunctions)\\
10 molecules, H$_2$--CO$_2$ & $10^{-9}$ m & $T/|V|$ & $1/2$ at equilibrium geometry\\
Thermal equilibrium & all & energy per quadratic d.o.f. & $k_BT/2$\\
White dwarfs & $10^{7}$ m & Chandrasekhar mass & $1.44\,M_\odot$, observed\\
The Sun & $10^{9}$ m & $K/|U|$ & $\approx1/2$ (to $\sim10\,\%$, uniform-sphere estimate)\\
Galaxy clusters (Coma and $\sim20$ others) & $10^{23}$ m & $K/|U|$ & $\approx1/2$ (to $\sim20\,\%$, projection)\\
Cosmological coupling & $10^{26}$ m & $\beta_m/\Omega_m$ & $1/2$ by prediction; consistent with Planck in 17 chains\\\hline
\end{tabular}
\end{table}

For converged Hartree--Fock wavefunctions the ratio is exact by construction: any wavefunction optimised under a uniform scaling of the coordinates satisfies
the quantum virial theorem. The atomic rows therefore show that quantum mechanics obeys the theorem; they are a statement of the theorem's reach, not an
independent measurement. The gravitational rows are measurements, at the precision the systems allow. The span from the atom to the cluster is 33 orders of
magnitude in size; with the cosmological coupling it reaches the Hubble scale.

\section{What the kinetic half is: the thermodynamic identity}
Clausius derived the theorem from Newton's second law. The derivation describes the bound state; it says nothing about the transition into it. Consider a
system that falls from rest at infinity ($E_i=0$) into a stable bound state with $E_f=\langle K\rangle+\langle V\rangle<0$, releasing the difference to its
surroundings.
\begin{enumerate}
\item First law: the heat released is $Q=E_i-E_f=-(\langle K\rangle+\langle V\rangle)>0$.
\item Second law: the transition is irreversible and produces entropy $\Delta S=Q/T$.
\item Landauer: the minimum cost of producing $\Delta S$ irreversibly is $E_{\rm L}=T\Delta S=Q$. The temperature cancels.
\item The $1/r$ equilibrium condition $2\langle K\rangle+\langle V\rangle=0$ gives $-E_f=\langle K\rangle$.
\end{enumerate}
Together,
\begin{equation}
\langle K\rangle=Q=T\Delta S=E_{\rm L}=\tfrac12|\langle V\rangle| .
\end{equation}
The hydrogen atom shows it directly: an electron captured from rest into the ground state emits a 13.6\,eV photon, the binding energy, which equals the
electron's mean kinetic energy in that state; the potential energy is $-27.2$\,eV. The energy carried away when a $1/r$ system binds equals its kinetic half.

This does not derive the virial theorem from thermodynamics; the fourth step is the mechanical condition. What it does is identify the kinetic half with a
definite thermodynamic quantity: the irreversible cost of the binding transition, carried off to the surroundings and recorded there as entropy at
$k_BT\ln2$ per bit. Jacobson (1995) obtained the Einstein equations from $\delta Q=T\,\delta S$ on local Rindler horizons; the same relation applied at the
boundary of a forming bound state returns the virial partition. The identity holds wherever the binding energy is released --- radiated by an atom, by a
contracting gas cloud, by a forming star. For collisionless dark matter, which relaxes without radiating, the energy is redistributed rather than released,
and the identity applies to the information written in the relaxation, not to emitted heat; this is the step IAM takes at cosmological scale.

\section{From the theorem to the cosmological coupling}
IAM adds to the horizon first law of Jacobson and Cai--Kim the entropy of the classical records produced as matter decoheres gravitationally,
\begin{equation}
-dE=T_H\,d\!\left(S_{\rm geo}+S_{\rm info}\right).
\end{equation}
The share of the matter density that enters the record term is the kinetic half of the binding energy of collapsed structure. The virial theorem fixes it:
\begin{equation}
\beta_m=\frac{\Omega_m}{2}=0.1577\qquad(\Omega_m=0.3153,\ \text{Planck 2018}).
\end{equation}
The coupling is a prediction fixed by the measured matter density; it enters every chain at this value and is never sampled.

It is useful to see where the $1/2$ sits in a real halo population. Writing $\beta_m=\Omega_m\,f_{\rm coll}\,\eta_{\rm vir}$, with $f_{\rm coll}$ the
fraction of matter in collapsed halos and $\eta_{\rm vir}$ the share of their binding energy in random, relaxed motion, the theorem requires
$f_{\rm coll}\eta_{\rm vir}=1/2$. Real halos are never perfectly relaxed: they are still accreting and merging (unrelaxed halos sit further from
equilibrium than relaxed ones), they carry ordered rotation, and some of their support is non-thermal (feedback, magnetic fields, cosmic rays). Each moves the
measured ratio away from the equilibrium value. In simulations, halos are conventionally called relaxed when $2T/|U|<1.35$ (Neto et al.\ 2007), and the
mass-weighted population sits somewhat above 1. The decomposition is therefore a description of how the equilibrium value is approached in a population,
not a separate confirmation of it; $\beta_m=\Omega_m/2$ rests on the theorem.
\authorhold{The source papers tabulate six N-body ``virial efficiencies'' (0.76--0.90) and six ``effective exponents'' (2.5--3.8). The published papers report
$2T/|U|$ of about 1.05--1.4, so 0.76--0.90 matches the reciprocal $|U|/2T$; the exponent values are not reported as such in the cited papers
(\texttt{verification/virial/NBODY\_TRACE.md}). The tables are not printed until each entry is traced to a figure or table in its source.}

\section{What the coupling predicts}
\subsection{The activation function and the growth coupling}
The record accumulates as structure forms. Integrating the record rate through the horizon first law gives the activation function
\begin{equation}
E(a)=\exp\!\left(1-\frac1a\right)=e^{-z},
\end{equation}
with $E\to0$ at early times ($E(z{=}10)=4.5\times10^{-5}$), $E(1)=1$ today, $dE/da>0$ throughout, and $E\to e$ as $a\to\infty$. It reaches 10, 50 and
90\,\% of today's value at $z=2.30$, $0.69$ and $0.11$. Matter perturbations grow with the effective coupling
\begin{equation}
\ddot\delta_m+2H\dot\delta_m-\tfrac32\Omega_mH_0^2a^{-3}\mu(a)\,\delta_m=0,\qquad
\mu(a)=\frac{H^2_{\Lambda{\rm CDM}}(a)}{H^2_{\Lambda{\rm CDM}}(a)+\beta_mE(a)H_0^2},
\end{equation}
while light is unaffected, $\Sigma=1$: photons on null worldlines accumulate no proper time and take no part in decoherence.

\begin{center}\small
\begin{tabular}{cccc}
\hline
$z$ & $E(a)$ & $\mu$ & growth coupling reduced by\\\hline
0 & 1.000 & 0.864 & 13.6\,\%\\
0.3 & 0.741 & 0.922 & 7.8\,\%\\
0.5 & 0.607 & 0.948 & 5.2\,\%\\
0.7 & 0.497 & 0.966 & 3.4\,\%\\
1.0 & 0.368 & 0.982 & 1.8\,\%\\\hline
\end{tabular}
\end{center}

\subsection{Amplitude of structure}
In the Planck chains $\sigma_8$ moves from $0.809$ to $0.800$ (Level 2; from $0.814$ to $0.801$ at Level 1), with every standard parameter within $0.1\sigma$
of its $\Lambda$CDM value: the lower amplitude is growth physics, not a parameter shift. The Level 2 chain gives $S_8=0.822\pm0.011$, in line with the
KiDS-Legacy cosmic-shear value $S_8=0.815^{+0.016}_{-0.021}$.

\subsection{Two expansion rates}
Photon-sector probes (CMB acoustic scale, BAO angles) read the $\Lambda$CDM rate; matter-sector probes calibrated on timelike worldlines read
\begin{equation}
H_0^{(m)}=H_0^{(\gamma)}\sqrt{1+\beta_m}=67.16\times1.0759=72.26\ \text{km\,s}^{-1}\text{Mpc}^{-1}
\end{equation}
(Level 2 chain; with Planck's $67.36$, $72.48$). SH0ES measures $73.04\pm1.04$ ($0.75\sigma$); Planck $67.36\pm0.54$ against the chain's $67.16\pm0.47$
($0.37\sigma$). Supernova distances follow $\Lambda$CDM geometry; their $H_0\approx73$ comes from the Cepheid calibration (Chapter~\ref{ch:dsvalidation}).

\subsection{Growth and geometry give different $\Omega_m$}
A growth measurement interpreted in $\Lambda$CDM infers $\Omega_m^{\rm growth}(z)=\Omega_m\,\mu(z)$: $0.299$ at $z=0.5$, $5.2\,\%$ below the geometric
value. The DESI DR1 full-shape plus BAO analysis gives $\Omega_m=0.2962\pm0.0095$ (effective redshift near 0.5), $0.3\sigma$ from this value. The DESI
analysis uses a $\Lambda$CDM template; a template built from the modified growth is the clean version of this test.

\subsection{The $E_G$ statistic}
$E_G$ compares galaxy--lensing cross-correlation with galaxy clustering and velocities. With $\Sigma=1$ the lensing kernel is unchanged and the growth rate is
lower, so $E_G=\Omega_m/f_{\rm IAM}(z)$ lies above the general-relativity value at every redshift. All the input data are published; the comparison with the
IAM growth curve is a direct near-term test.

\subsection{The virial ratio through cosmic history}
The ratio of the matter density to the record term, $R(a)=\Omega_ma^{-3}/[\beta_mE(a)]$, is $399$ at $z=2$, $20$ at $z=0.7$, $6$ at $z=0.3$, and $2$
today --- the last by construction, since $\beta_m=\Omega_m/2$ and $E(1)=1$. Matter equals the vacuum plus record term at $z=0.361$ ($\Lambda$CDM:
$0.295$). Today the record term is $18.7\,\%$ of $\Omega_\Lambda+\beta_m$; at $z=0.3$, $0.7$ and $1.5$ it is $14.6$, $10.3$ and $4.9\,\%$. The reading of
$R=2$ as the solution of the coincidence problem is developed in Part~5.

\section{Tests}
\begin{itemize}
\item \emph{Euclid} measures growth (spectroscopy) and lensing (imaging) on the same volume. The prediction is $\mu_0-1=-0.136$ with $\Sigma_0=0$; the
full survey's $\sigma(\mu_0)\approx0.04$ gives about $3.4\sigma$. The clustering and weak-lensing products of DR1 are expected in mid-2027.
\item \emph{DESI Year 5} $f\sigma_8(z)$ at $\sim1\,\%$ per bin: a suppression that follows $E(a)$ (largest at low $z$, $1.8\,\%$ at $z=1$), not a constant
rescaling.
\item $E_G$ from published galaxy--lensing data.
\item What would count against the virial coupling: $\mu_0$ consistent with zero at Euclid precision; correlated $\mu\neq1$, $\Sigma\neq1$; a
redshift-independent growth suppression; or, with $\Omega_m$ revised, growth data that still require the old $\beta_m$.
\end{itemize}

\authorhold{For discussion before inclusion (speculative or interpretive; Part 5 or out): the identification of dark matter and dark energy as the potential
and kinetic halves, and ``what is spacetime made of'' (paper 3, \S3, \S6); the arrow-of-time sections and the England, Rovelli and Smolin discussions
(paper 4, \S3.3, \S6); the black-hole ``encoding vault'' and the $\sigma_{\rm crit}\approx4$\,km\,s$^{-1}$ satellite threshold (paper 4, \S4); the
three-channel split of $\beta_m$ and the three-way cluster ratio test (paper 2, \S3.3, \S6.2); $w_{\rm info}=-4/3$ and $z_t=0.718$, which need a modified
background (paper 3, \S4.2); the DESI phantom crossing as a fitting artifact (paper 3, \S4.3; open until the real-data two-ruler test is run).}
